\section{Introduction}
\label{sec:intro}

Voice assistants built on full-duplex speech models
\citep{defossez2024moshi,zeng2024glm4voice,xu2025qwen25omni} face a
structural tension. The conversational loop---listening while speaking,
turn-taking inside a $\sim$200--300\,ms budget
\citep{stivers2009universals}---demands a small, low-latency model close to
the user. Answer quality on knowledge- and reasoning-heavy queries demands a
large one. The natural architecture is a \emph{talker--thinker} split,
echoing the System~1/System~2 distinction \citep{kahneman2011thinking}: a
small duplex ``talker'' holds the live session and answers by default, and a
large cloud ``thinker'' is consulted only when needed. This paper is about
the load-bearing decision in that architecture: \textbf{when should the
talker hand off?}

We deliberately restrict ourselves to \emph{zero-training} gate signals ---
quantities that can be read off the small model as-is, without fine-tuning
it: decode-time scalars (entropy, margin), linear probes on hidden states
\citep{alain2017understanding}, and verbalized self-evaluation, i.e.\
$\ptrue$-style ``Would you answer this correctly?'' prompts
\citep{kadavath2022language}. The constraint is practical (duplex fine-tunes
are expensive and fragile to retrain) but it also sharpens the scientific
question: \emph{does a small full-duplex model already know what it does not
know, and can that knowledge be read out in time to act on it?}

Our answer, developed over a controlled study of ten
backbones/fine-tunes and both input modalities, is: \textbf{the knowledge is
there; the readouts are what break, and they break in systematic,
duplex-fine-tune-induced ways.} Concretely:

\begin{enumerate}
\item \textbf{Standard probes flatter to deceive (\S\ref{sec:rq1}).} A
  final-layer hidden-state probe reaches AUC 0.822, but an audit shows it is
  mostly a query-\emph{type} recognizer: a pool-identity oracle alone scores
  0.715, and leave-one-pool-out (LOPO) transfer \emph{inverts} on math
  (0.372). Pre-answer $\ptrue$ needs no training, beats the probe on most
  backbones, and scores 100\%-fail SimpleQA traps at 0.945 --- but only if
  asked \emph{before} the model answers; after answering, it endorses its own
  confident-wrong answer (0.945 $\rightarrow$ 0.360).
\item \textbf{Duplex fine-tuning overwrites readouts, not knowledge
  (\S\ref{sec:duplex}).} Across matched pairs (raw vs.\ omni-streaming vs.\
  duplex fine-tunes of the same backbone), LOPO probe transfer degrades in
  the order raw $>$ omni-streaming $>$ duplex, with label coverage
  controlled. A full layer$\times$position sweep localizes the damage: on
  MiniCPM-o~4.5, LOPO math is 0.931 at layer 22 but 0.366 at layer 35, with
  a sharp cliff over the last four layers, only at the last-token position
  --- exactly where a streaming head must encode turn-control state.
\item \textbf{Each readout has one blind modality (\S\ref{sec:speech}).}
  The late-layer cliff is text-input-specific: under audio input the final
  layer does not invert (0.936). Conversely, verbalized introspection
  collapses on audio input (trap $p_{\mathrm{yes}}$ $0.055 \rightarrow
  0.556$) on both MiniCPM duplex generations while omni-streaming and
  frozen-backbone controls stay intact. Three controlled arms (ASR audit,
  filler audio, audio+text duplication) pin the mechanism to \emph{binding}:
  the verbal self-check runs over text-token pathways that audio embeddings
  do not feed, while a probe reads the same information at 0.93+.
\item \textbf{End-to-end training creates the shared core
  (\S\ref{sec:speech}).} Text$\leftrightarrow$audio probe transfer is
  $\sim$0.86 on every end-to-end model but dead ($0.52$--$0.60$) on
  Freeze-Omni's frozen backbone with identical text weights --- so the
  practical recipe ``calibrate on cheap text, deploy on speech'' is valid
  for end-to-end models only.
\item \textbf{The gate wins as a system (\S\ref{sec:system}).} The mid-layer
  probe decides in 20\,ms (text) / 45\,ms (audio), before the first output
  token. Against the model's own \texttt{enable\_thinking} mode, gated-cloud
  escalation dominates on both accuracy and latency at every escalation
  rate, because the gate flags knowledge/trap failures that thinking cannot
  fix --- external escalation is \emph{necessary}, not merely better.
\end{enumerate}

\begin{figure}[t]
  \centering
  \includegraphics[width=\linewidth]{concept_probe_vs_ppl}
  \caption{Conceptual contrast between a supervised probe in hidden-state
  space and what output-distribution scalars (entropy/perplexity) can see.
  \todonote{regenerate this figure with English labels (current version has
  Chinese axis/legend text); source:
  \texttt{figures/concept\_probe\_vs\_ppl.py}}}
  \label{fig:concept}
\end{figure}

This is explicitly \textbf{step~1 of a two-step design}. Step~2 --- how the
thinker's result is returned into a live duplex session (injection protocol,
authority framing, stall phrases) --- is scoped but not solved here; our
streaming experiments already surface its first real obstacle (the talker
pushes back on injected results that conflict with its own reasoning,
\S\ref{sec:discussion}), and our overlap analysis (\S\ref{sec:system})
quantifies its timing budget.
